\section{Length of a Loop}
\label{sec:ll-loop-length}

\problemstatement{If a linked list contains a cycle, return the number of
nodes in the cycle; otherwise return $0$.}

\begin{patternbox}
Once Floyd's pointers meet inside the loop, \textbf{keep one pointer fixed
at the meeting node and walk the other around until it returns} -- the
number of steps is the loop length. The collision guarantees both pointers
are on the cycle.
\end{patternbox}

\subsection*{Walk once around from the meeting point}

Run the tortoise--hare until \texttt{slow} and \texttt{fast} coincide; that
node lies on the cycle. Now advance a single pointer from that node,
counting steps, until it comes back to the same node. Since it started on
the cycle, it traverses exactly the cycle once, and the count is the loop's
length. If the pointers never meet (fast reaches null), the length is $0$.

\begin{ideabox}[A Full Lap from Any Cycle Node Measures It]
Any node known to be on the cycle can measure the cycle's size: leave it in
place and count steps until a moving pointer returns to it. The Floyd
meeting point conveniently provides such a node.
\end{ideabox}

\subsection*{Algorithm}

\begin{enumerate}
  \item Run slow/fast until they meet; if no meeting, return $0$.
  \item From the meeting node, advance a pointer counting steps until it
        returns to the meeting node.
  \item Return the count.
\end{enumerate}

\subsection*{Dry run}

\dryrun{$1\to2\to3\to4\to2$ (loop $2\to3\to4\to2$, length $3$).}

\begin{itemize}
  \item Pointers meet inside $\{2,3,4\}$. From the meeting node, count a
        lap: $3$ steps back to it.
  \item Length $=3$.
\end{itemize}

\subsection*{C++ implementation}

\begin{lstlisting}
#include <bits/stdc++.h>
using namespace std;

struct ListNode {
    int val; ListNode* next;
    ListNode(int v) : val(v), next(nullptr) {}
};

int loopLength(ListNode* head) {
    ListNode* slow = head;
    ListNode* fast = head;
    while (fast && fast->next) {
        slow = slow->next;
        fast = fast->next->next;
        if (slow == fast) {                       // on the cycle
            int len = 1;
            ListNode* p = slow->next;
            while (p != slow) { p = p->next; len++; }
            return len;
        }
    }
    return 0;
}

int main() {
    ListNode* a = new ListNode(1);
    a->next = new ListNode(2);
    a->next->next = new ListNode(3);
    a->next->next->next = new ListNode(4);
    a->next->next->next->next = a->next;          // 4 -> 2
    cout << loopLength(a) << "\n";                 // 3
    return 0;
}
\end{lstlisting}

\begin{complexitybox}
\textbf{Time Complexity.} $O(n)$. \textbf{Space Complexity.} $O(1)$.
\end{complexitybox}

\begin{takeawaybox}
Loop length is one lap counted from the tortoise--hare meeting node. Cycle
detection, its start, and its length are three questions answered by the
same slow/fast collision -- keep the trio together in memory.
\end{takeawaybox}
